%=================================================================
% PLANTILLA -- Respuesta a los revisores -- Segundo Corte
% LabDST RoboCoffee, ITESO.
%
% Como usarla:
%   1. Copien esta carpeta junto a la carpeta de su paper revisado
%      (las dos carpetas deben ser HERMANAS) y cambien "N" por el
%      numero de su equipo, en el nombre de la carpeta y del .tex.
%   2. Ajusten la ruta de \LabDSTpaper: relativa, SIN extension,
%      entre comillas y en UNA sola linea.
%   3. En el paper, marquen cada cambio con \linelabel{ln:...} (texto)
%      o \label{sec:/fig:/tab:/eq:...} (secciones, figuras, tablas,
%      ecuaciones).
%   4. Escriban un \LabDSTcomment por CADA clave del PDF de
%      observaciones (R1.1, R1.2, ..., R4.n), en el mismo orden.
%   5. Compilen primero el paper (pdflatex -> bibtex -> pdflatex ->
%      pdflatex) y despues esta carta DOS veces (pdflatex x2).
%      Si aparece "??", revisen la ruta o el nombre de la etiqueta.
%
% La carta se escribe en ingles, como el paper. El comentario del
% revisor se copia TEXTUAL desde el PDF de observaciones.
%
% \LabDSTcomment{Clave}{Estado}{Ubicacion}{Comentario}{Respuesta}
%   Estado:    A = Addressed
%              P = Partially addressed
%              N = Not addressed (justified)
%   Ubicacion: \rln{x} -> "line NN"     \rlns{x}{y} -> "lines NN--MM"
%              \Rln, \Rlns (inicio de oracion), Section~\ref{sec:...},
%              Fig.~\ref{fig:...}, Table~\ref{tab:...}, \eqref{eq:...}
%   \LabDSTquote{...} cita el texto NUEVO del paper.
%
% Borren los comentarios <...> y todo lo que este entre < >.
%=================================================================
\documentclass{LabDSTresponse}
\LabDSTpaper{"../Equipo 3 - tercer corte/Equipo 3 - tercer corte"}

\begin{document}
	\LabDSTtitle{Second Cut}{Team 3}{Servo Characterization for the SOFIA Modular Kinetic
		Surface}
	
	%-----------------------------------------------------------------
	\LabDSTsection{Acknowledgements}
	
	The authors thank the reviewers for their careful evaluation of the first-cut manuscript. Their comments helped clarify the research scope, experimental methodology, and electrical compatibility. All comments were considered; the responses below follow the keys of the review report (R1.1, R1.2, \ldots). Line, section, figure, table, and equation numbers refer to the revised manuscript that accompanies this letter.
	
	\textbf{Main changes.} (i) Refined problem statement and objectives (Section~\ref{sec:intro}). (ii) Expanded characterization methodology (Section~\ref{sec:methods}). (iii) Documented RS-485 voltage-level adaptation (Section~\ref{sec:rs485_logic}).
	
	\LabDSTsummary
	
	%-----------------------------------------------------------------
	\LabDSTreviewer{1}{Format, LaTeX and template compliance}
	
	\LabDSTcomment{R1.1}{A}{}
	{Discussion y Conclusions son solo encabezados, y no existen Author Contributions, Data Availability (con la URL del
		repositorio), Acknowledgment (con la declaración de uso de IA) ni Conflicts of Interest. Era la observación R1.5 del primer
		corte.}
	{}
	
	\LabDSTcomment{R1.2}{A}{}
	{Los nombres de los gestos están en español: INICIO, FIN, RECARGAR, ESCUDO y ATAQUE. También hay un error de
		espacio: “in thisstudy” (Línea 162).
	}
	{}
	
	\LabDSTcomment{R1.3}{A}{}
	{La Fig. 1 mide 6 cm de ancho y su texto interno (“CH 1”, “JST-PH (MALE)”) no se lee impreso. Tampoco muestra el riel
		de 5 V, la fuente, el transceptor RS-485 ni el dispositivo de corona. El esquemático (Fig. 4) es ilegible a una columna, y la
		PCB (Fig. 5) no tiene cotas ni identificación de pistas.}
	{}
	
	\LabDSTcomment{R1.4}{A}{}
	{La lista de autores no incluye a la profesora Graciela Baca, y la afiliación no indica el departamento.}
	{}
	
	\LabDSTcomment{R1.5}{A}{}
	{El BOM conserva elementos sin describir (“Pro3”, “Axis 1 (9 × 9 mm)”, “Inser M3x5x4”) y una fila “MG90S / SG90”,
		aunque la Tabla IX descarta el SG90. El PCA9685 se atribuye a “NXP Semi”, pero se usa como módulo comercial. Faltan
		el divisor del RS-485 como componente de la PCB y la fuente de alimentación}
	{}
	
	%-----------------------------------------------------------------
	\LabDSTreviewer{2}{Scientific Writing and Language}
	
	\LabDSTcomment{R2.1}{A}{}
	{El abstract describe el método propuesto, pero no da ningún resultado, aunque ya existen: la medición de 5.00 V y el divisor,
		y las métricas de los gestos. Además, describe el dispositivo de corona como si fuera parte del trabajo, cuando no se ha
		diseñado.}
	{}
	
	\LabDSTcomment{R2.1}{A}{}
	{Hay afirmaciones que los datos no sostienen:
		• “power supply stability and protection against voltage sags must be managed through software-level trajectory
		smoothing” (Línea 424): el suavizado reduce los picos, pero no sustituye una fuente y una distribución de potencia dimensionadas;
		• “little mechanical chatter and motion desynchronization occurred, establishing 350◦
		/s as the maximum operational
		threshold” (Línea 435): no hay tabla, gráfica ni criterio que defina el umbral;
		• “Preliminary observations from the velocity sweeps confirm small timing variations between identical micro-servos”
		(Línea 445): no se reporta ninguna medición de tiempo.}
	{}
	
	%-----------------------------------------------------------------
	\LabDSTreviewer{3}{Technical Content: Materials and Methods, and Design and Implementation}
	
	\LabDSTcomment{R3.1}{A}{}
	{El esquemático tiene tres errores eléctricos que el texto no menciona:
		• Lógica del PCA9685 a 5 V. La terminal VCC de J1 y J2 está en la red VCC, que es el pin de 5 V de la ESP32-C3.
		El pin de 3.3 V no está conectado. Así, los PCA9685 operan su lógica a 5 V y, en los módulos comerciales, las
		resistencias de pull-up de SDA y SCL llevan 5 V a GPIO0 y GPIO1, que no toleran 5 V. La Tabla V dice lo contrario:
		“Uses 3.3 V logic power supply”.
		4
		• LED en pines de arranque. El LED RGB está en GPIO8 (azul) y GPIO9 (verde), y OE en GPIO2: son los tres pines
		de arranque que la Tabla IV pide verificar. Un LED con su resistencia a tierra puede bajar GPIO9 en el arranque y
		dejar la ESP32-C3 en modo de descarga. Además, R1–R3 no tienen valor en el esquemático.
		• Potencia de los servos por JST-PH. El riel V+ de los servos solo entra por J3/J4, conectores JST-PH de 2.0 mm
		que soportan alrededor de 2 A por contacto. El propio artículo estima más de 22 A para 32 servos (Línea 411). La
		arquitectura de SOFIA usa JST-VH (observación R3.6 del primer corte).
		Tampoco hay capacitores de desacoplo ni de bulk}
	{}
	
	\LabDSTcomment{R3.2}{A}{}
	{La sección de potencia usa “approximately 250 mA during movement” y “peak stall currents reaching up to 700 mA” sin
		fuente, mientras la Tabla VI dice que la corriente de bloqueo “Not reported” requiere medición. No se dimensiona la fuente
		ni se dice qué corriente soporta la PCB.}
	{}
	
	\LabDSTcomment{R3.3}{A}{}
	{El dispositivo de corona con interruptores de límite, que es la base del diferenciador, solo se menciona: no hay CAD, figura,
		número ni posición de los interruptores, resolución angular, ni cómo se registra el tiempo (interrupciones de la ESP32-C3,
		timer, frecuencia). Además, el protocolo (Línea 244) habla de medir el “transient velocity profile”, pero unos interruptores
		en posiciones fijas dan tiempos de paso, no un perfil continuo. No se explica cómo se ajusta $\tau_m$m a esos tiempos.}
	{}
	
	\LabDSTcomment{R3.4}{A}{}
	{No se describe la arquitectura del firmware de la ESP32-C3 ni el formato de los mensajes entre la PC y la ESP32-C3 (Wi-Fi
		o Bluetooth), ni el uso del RS-485 entre placas (Línea 180). El valor de $\alpha_m$ del Algoritmo 1 no se da.}
	{}
	
	%-----------------------------------------------------------------
	\LabDSTreviewer{4}{D-Lab Methodology, Results, and Validation}
	
	\LabDSTcomment{R4.1}{P}{}
	{La pregunta de investigación es si estimar $\tau_m$m e integrarla al controlador mejora la sincronización. Ninguna de las dos partes
		avanzó: no hay dispositivo ni $\tau_m$m medido, y el resultado que se reporta es el reconocimiento de gestos, que no responde la
		pregunta. La Línea 445 pospone la medición al siguiente corte.}
	{}
	
	\LabDSTcomment{R4.2}{P}{}
	{Las métricas de la Tabla XIV no son consistentes:
		• La Línea 490 dice que ESCUDO se confundió con RECARGAR. Pero RECARGAR tiene precisión de 100%, lo que
		no es posible si hubo ESCUDOs detectados como RECARGAR.
		• FIN tiene precisión de 70%, es decir, tres detecciones falsas de FIN que no aparecen en ninguna otra fila.
		• La columna Accuracy es igual a Recall en todas las filas: por clase, eso es la tasa de aciertos, no la exactitud.
		• La Tabla XV tiene 4 intentos de ESCUDO y la Tabla XIV 12, sin decir si son la misma prueba. La Línea 499 aclara
		que las mediciones preliminares son de una versión anterior, pero no dice si la Tabla XIV lo es.
		Con 6–16 intentos por clase y sin intervalos de confianza, los porcentajes tienen mucha incertidumbre}
	{}
	
	\LabDSTcomment{R4.3}{P}{}
	{La pregunta y el objetivo 3 hablan de “a 32-actuator kinetic surface” y de igualar los tiempos de asentamiento de todos
		los actuadores, pero el alcance declarado es de dos servos de un módulo (Línea 162). El objetivo 1 (diseñar la PCB) no se
		relaciona con la pregunta, y los gestos no se relacionan con ningún objetivo.}
	{}
	
	\LabDSTcomment{R4.4}{P}{}
	{Sigue sin haber análisis Discover del demostrador: la Línea 113 habla de “the functional decomposition of the demonstrator”, pero no se muestra. No se declaran el repositorio en GitHub ni la bitácora, y no hay incrementos de integración
		documentados.}
	{}
	
	
\end{document}
